\section{Exact computational certificates}
\label{app:computational-certificates}

The computer-assisted statements in this volume use exact integer or
finite-field arithmetic.  Floating-point computations are never used as
proof of a Galois-group, separability, or valuation assertion.  The complete
scripts and captured outputs are shipped in the source archive under
\nolinkurl{computation/}.

\subsection{All-order rectangle}

The command
\begin{center}
\texttt{python3 computation/trace\_allq\_modp\_audit.py
  --bound 1001 --candidate-count 10}
\end{center}
examines, for every eligible odd degree, at most the ten largest
strict-window primes in descending order.  It checks every relevant order
residue class, records the prime-degree fallback rows, and emits the combined
certificates at \(m=27\) and \(m=989\).  The seven \(m=27\) discriminant
rows use the values displayed in the main text; the search does not replace
them by resultants.

\subsection{Fixed-order polylogarithmic window}

The command
\begin{center}
\texttt{python3 computation/polylog\_window\_prime\_audit.py
  --bound 1001}
\end{center}
uses exact arithmetic in $\F_p$ to test every strict-window prime
$(m+1)/2<p<m$ for each $\tau\in\{1,2,3,4\}$ and odd
$7\le m\le1001$.  It constructs the residual polynomial $R_r$ from
\eqref{eq:R-general} and checks $\gcd(R_r,R_r')=1$.  There are no window
exceptions for $\tau=1,2,3$; for $\tau=4$ the sole window exception is
$m=17$, where the unconditional prime-degree part of
\Cref{thm:polylog-global} proves $S_{17}$.  Thus the exception concerns only
the window certificate, not the conclusion of \Cref{cor:polylog-window}.

The frozen files have SHA-256 hashes
\begin{itemize}
\item \nolinkurl{polylog_window_prime_audit.py}:\\
\texttt{\small\seqsplit{dd42b574f3eb5144cd49316530911f4cee0e3722ed03dda7382bd10ec8a31cd2}};
\item \nolinkurl{polylog_window_prime_audit_1001.txt}:\\
\texttt{\small\seqsplit{8f952488d2011c1391ae15a4f607715c7fc7de73a00ce8cae00fa973f938cbdf}}.
\end{itemize}

\subsection{Original tower through one million}

The command
\begin{center}
\texttt{python3 computation/window\_prime\_selection\_audit.py
  --bound 1000000}
\end{center}
checks the original \(q=1\) tower.  For each odd composite degree it first
tests the largest strict-window prime and then the next candidate if needed.
The output records the two failures of the largest-prime choice and the
successful second-prime certificates.

\subsection{Morse and derivative certificates}

The lower-degree exact Morse certificates are generated by
\begin{center}
\texttt{python3 computation/trace\_morse\_certificates.py
  --max-m 55 --prime-bound 3000}.
\end{center}
The residual even-derivative computations are implemented independently in
the supplied C++ source and their captured outputs are included for bounds
\(1000\) and \(5000\).  These finite certificates support exactly the ranges
stated in the theorem ledger; they are not extrapolated to a universal DS
theorem.

\subsection{Reproducibility convention}

Each source-archive file is covered by the top-level SHA-256 manifest.  A
successful independent reproduction consists of:

\begin{enumerate}
\item verifying the manifest;
\item compiling the monograph with the documented pdfLaTeX command;
\item running the quick certificate suite;
\item optionally running the extended \(10^6\)-degree suite and comparing
its terminal summary with the captured output.
\end{enumerate}

The source archive records the interpreter/compiler versions used for the
frozen build.  No network access is required to compile the PDF or rerun the
certificates.
